\section*{Chapter \ref{ch:g12:stereochemistry} --- Stereochemistry: Chirality and Isomers}

\begin{solution}{exo:g12:stereochemistry:1}
Butan-2-ol: carbon 2 (\ce{H}, \ce{OH}, \ce{CH3}, \ce{CH2CH3}).
Propan-2-ol: none (carbon 2 carries two identical \ce{CH3}).
2-Chlorobutane: carbon 2. Lactic acid: carbon 2 (\ce{H}, \ce{OH},
\ce{CH3}, \ce{COOH}).
\end{solution}

\begin{solution}{exo:g12:stereochemistry:2}
First: the two \ce{Cl} (priority over \ce{H}) on the same side: \omterm{def:g12:stereochemistry:z-e}{Z}.
Second: on the left, \ce{CH3} has priority, upper side; on the right,
\ce{CH2CH3} has priority, lower side: \omterm{def:g12:stereochemistry:z-e}{E}.
\end{solution}

\begin{solution}{exo:g12:stereochemistry:3}
\omterm{def:g12:stereochemistry:chiral}{Chiral}: the glove, the screw. Not \omterm{def:g12:stereochemistry:chiral}{chiral}: the spoon, the cube (each has
a mirror plane). \ce{CHFClBr} is \omterm{def:g12:stereochemistry:chiral}{chiral} (four different \omterm{def:g7:atoms-and-molecules:atom}{atoms} on carbon);
\ce{CH2Cl2} is not.
\end{solution}

\begin{solution}{exo:g12:stereochemistry:4}
A \omterm{def:g6:pure-substances-and-mixtures:mixture}{mixture} of equal amounts of two \omterm{def:g12:stereochemistry:enantiomer}{enantiomers}. Racemic carvone contains
both \omterm{def:g12:stereochemistry:enantiomer}{enantiomers}, so the nose receives both smells at once.
\end{solution}

\begin{solution}{exo:g12:stereochemistry:5}
No: carbon 1 carries two identical \omterm{def:g7:atoms-and-molecules:atom}{atoms}, \ce{H} and \ce{H}; exchanging
them gives the same \omterm{def:g7:atoms-and-molecules:molecule}{molecule}.
\end{solution}

\begin{solution}{exo:g12:stereochemistry:6}
\setchemfig{atom sep=2em}
\chemfig{C(-[2]CH_2CH_3)(-[4]H_3C)(<[:-30]OH)(<:[:-150]H)} \quad mirror
\quad \chemfig{C(-[2]CH_2CH_3)(-[0]CH_3)(<[:-150]HO)(<:[:-30]H)}
\end{solution}

\begin{solution}{exo:g12:stereochemistry:7}
2-Chlorobutane: 2 (one \omterm{def:g12:stereochemistry:chiral}{asymmetric carbon}). 2,3-Dichlorobutane: 3 (two
\omterm{def:g12:stereochemistry:chiral}{asymmetric carbons} with alike halves: a pair of \omterm{def:g12:stereochemistry:enantiomer}{enantiomers} and a meso
form). Pent-2-ene: 2 (\omterm{def:g12:stereochemistry:z-e}{Z} and \omterm{def:g12:stereochemistry:z-e}{E}).
\end{solution}

\begin{solution}{exo:g12:stereochemistry:8}
\omterm{def:g12:stereochemistry:enantiomer}{Enantiomers}; \omterm{def:g12:stereochemistry:diastereomer}{diastereomers}; \omterm{def:g12:stereochemistry:diastereomer}{diastereomers}; the same \omterm{def:g7:atoms-and-molecules:molecule}{molecule} (meso
tartaric acid is superimposable on its mirror image).
\end{solution}

\begin{solution}{exo:g12:stereochemistry:9}
In the staggered \omterm{def:g12:stereochemistry:conformation}{conformation} with opposite \ce{CH3} groups, the large
groups are as far apart as possible: it is the most stable. In the
eclipsed one, the two \ce{CH3} face each other and crowd one another.
\end{solution}

\begin{solution}{exo:g12:stereochemistry:10}
Its two halves, each \ce{CH(OH)COOH}, are mirror images of each other: a
plane between the two central carbons cuts the \omterm{def:g7:atoms-and-molecules:molecule}{molecule} into two mirror
halves. A \omterm{def:g7:atoms-and-molecules:molecule}{molecule} with a mirror plane is its own mirror image: not
\omterm{def:g12:stereochemistry:chiral}{chiral}.
\end{solution}

\begin{solution}{exo:g12:stereochemistry:11}
\setchemfig{atom sep=1.7em}
(\omterm{def:g12:stereochemistry:z-e}{Z})-pent-2-ene \chemfig{C(-[:150]H_3C)(-[:210]H)=C(-[:30]CH_2-CH_3)(-[:-30]H)}
and (\omterm{def:g12:stereochemistry:z-e}{E})-pent-2-ene \chemfig{C(-[:150]H_3C)(-[:210]H)=C(-[:30]H)(-[:-30]CH_2-CH_3)}.
\end{solution}

\begin{solution}{exo:g12:stereochemistry:12}
As for tartaric acid, with \ce{Br} instead of \ce{OH} and \ce{CH3}
instead of \ce{COOH}: both \ce{Br} solid, both dashed (a pair of
\omterm{def:g12:stereochemistry:enantiomer}{enantiomers}), and one solid one dashed, which has a mirror plane: the
meso form. Three \omterm{def:g12:stereochemistry:stereoisomer}{stereoisomers}.
\end{solution}

\begin{solution}{exo:g12:stereochemistry:13}
Half. The chemist can separate the two \omterm{def:g12:stereochemistry:enantiomer}{enantiomers} (as Pasteur did, by
other means), or make only the useful enantiomer, with a reaction that
is itself \omterm{def:g12:stereochemistry:chiral}{chiral} (a \omterm{def:g12:stereochemistry:chiral}{chiral} catalyst, an enzyme, a \omterm{def:g12:stereochemistry:chiral}{chiral} starting
material).
\end{solution}

\begin{solution}{exo:g12:stereochemistry:14}
Three: (2E,4E), (2Z,4Z) and (2E,4Z), the last being the same \omterm{def:g7:atoms-and-molecules:molecule}{molecule} as
(2Z,4E) read from the other end.
\end{solution}

\begin{solution}{exo:g12:stereochemistry:15}
$2^3 = 8$ \omterm{def:g12:stereochemistry:stereoisomer}{stereoisomers}, that is 4 pairs of \omterm{def:g12:stereochemistry:enantiomer}{enantiomers}. Each
stereoisomer has one enantiomer and $8 - 2 = 6$ \omterm{def:g12:stereochemistry:diastereomer}{diastereomers}.
\end{solution}

\begin{solution}{pb:g12:stereochemistry:1}
\textbf{1.} \ce{C4H6O6}; $M = 4 \times 12.0 + 6 \times 1.0 + 6 \times 16.0
= \qty{150.0}{g/mol}$.

\textbf{2.} Two carboxyl groups (acid) and two hydroxyl groups
(\omterm{def:g11:functional-groups:alcohol}{alcohol}).

\textbf{3.} The two central carbons, each bonded to \ce{H}, \ce{OH},
\ce{COOH} and \ce{CH(OH)COOH}.

\textbf{4.} $2^2 = 4$.

\textbf{5.} Both with solid wedges.

\textbf{6.} Both \ce{OH} with dashed wedges; it cannot be superimposed on
L: it is D-tartaric acid, its enantiomer.

\textbf{7.} Turned about the central bond so that the two \ce{OH} face
the same way, the \omterm{def:g7:atoms-and-molecules:molecule}{molecule} has a plane between its two central carbons
that reflects one half onto the other.

\textbf{8.} No: it is meso-tartaric acid.

\textbf{9.} Of the four combinations (solid, solid), (dashed, dashed),
(solid, dashed), (dashed, solid), the last two are the same \omterm{def:g7:atoms-and-molecules:molecule}{molecule},
the meso form, because the two halves of the \omterm{def:g7:atoms-and-molecules:molecule}{molecule} are alike.

\textbf{10.} L and D: \omterm{def:g12:stereochemistry:enantiomer}{enantiomers}. L and meso: \omterm{def:g12:stereochemistry:diastereomer}{diastereomers}.

\textbf{11.} At \qty{169}{\celsius} too: \omterm{def:g12:stereochemistry:enantiomer}{enantiomers} have the same
physical properties.

\textbf{12.} No: it is a diastereomer of L, a different compound with
its own properties.

\textbf{13.} The receptors of the nose are \omterm{def:g12:stereochemistry:chiral}{chiral} and distinguish
mirror images; a thermometer, not \omterm{def:g12:stereochemistry:chiral}{chiral}, cannot.

\textbf{14.} No, a \omterm{def:g6:pure-substances-and-mixtures:mixture}{mixture} of two compounds; as a whole it is not
optically active, the two \omterm{def:g12:stereochemistry:enantiomer}{enantiomers} cancelling each other.

\textbf{15.} Only one enantiomer: each crystal was built of L or of D
alone.

\textbf{16.} \qty{5.0}{g} of each.

\textbf{17.} The meso form is a different compound, not part of the
\omterm{def:g12:stereochemistry:enantiomer}{racemic mixture} of L and D.

\textbf{18.} By dissolving each pile: the solutions turn polarised light
by the same amount in opposite directions (and the crystals are mirror
images).

\textbf{19.} Four: its two \omterm{def:g12:stereochemistry:chiral}{asymmetric carbons} carry different groups
(\ce{Cl} on one, \ce{OH} on the other), so no combination has a mirror
plane, and no two combinations are the same \omterm{def:g7:atoms-and-molecules:molecule}{molecule}.

\textbf{20.} Three: L, D and meso.
\end{solution}
